\section{Consent, and the Shape of the Gift}

The second option is consent: the will's agreement with what the intellect saw. Not resignation, which agrees the way a prisoner agrees with the wall, and not optimism, which agrees by not looking. Consent says: it is true that I do not hold myself in being, and I am willing that it be true.

\subsection{What consent is, and is not}

Three boundaries keep the word from doing more work than it can. Consent is not belief in anything further; a reader who has followed the argument and is unwilling to add revealed premises can perform it exactly as stated. It is not an emotion, though it may have emotional consequences; the vertigo was an act of intellect and its answer is an act of will, and neither is a mood. And it is not resignation to a fate. Resignation agrees with what cannot be helped, and its characteristic tone is the tone of a man who has stopped arguing with weather. Consent agrees with what is the case because it is the case, and the difference is visible in what each does next: resignation lowers its expectations, consent looks harder at the thing agreed to.

Which is why consent sounds passive and is not. Everything depends on what it consents \emph{to}, and to see that, one has to look again at the structure left standing at the end of section~6.

\subsection{The shape of a giving that cannot be repaid}

The source sustains everything and gains nothing. It cannot gain: no return current, no real relation back, no beatitude to be increased. Now ask what kind of giving that is---the question recoil never stays to ask.

Among creatures we know it only in traces, and the traces are instructive because of how we already judge them. The mother of a newborn who can know nothing of her, and will remember nothing of these months; the craftsman perfecting a joint no inspector will ever see; the anonymous donor; the man who does not correct a false attribution of his own work to someone else. In every case the absence of return is what makes us call the act generous rather than transactional. Giving that expects no return we call love, and we rank it higher precisely in proportion as the possibility of return recedes.

Now push the case to its limit. Here the absence of return is not an ethical achievement---not a return declined but a return metaphysically impossible---and the giving happens anyway: at every point of being, at every moment, to everything that is, and it could stop and does not. If the traces are called love because of the missing return current, what is to be said of the case where the current is missing by necessity?

\subsection{The judgment that God loves}

Aquinas turns this observation into a formal judgment---the summit of what unaided reason can say here---and the argument is short enough to follow in full. Everything that exists is good in the measure that it exists, ``since the existence of a thing is itself a good.'' God's will is the cause of all things, so anything that has existence or any other good has it as willed by God. ``To every existing thing, then, God wills some good. Hence, since to love anything is nothing else than to will good to that thing, it is manifest that God loves everything that exists.''\endnote{\emph{Summa theologiae} I, q.~20, a.~2, corpus (``God loves all existing things''); checked 2026-07-25, English at New Advent, Latin at Corpus Thomisticum. The definition doing the work---to love is to will good to something---is taken from Aristotle, \emph{Rhetoric} II, and is used throughout the \emph{Summa}'s treatment of love. The argument's premise that God's will is the cause of all things is established earlier, at I, q.~19, a.~4, and is not re-argued here.}

Then the difference from our loving, stated as a difference in the direction of causality: our love is awakened by a goodness already there, ``whereas the love of God infuses and creates goodness.''\endnote{Same corpus (\latin{amor Dei est infundens et creans bonitatem in rebus}); checked 2026-07-25 as above. Our love, Aquinas says, ``is not the cause of its goodness; but conversely its goodness, whether real or imaginary, calls forth our love.'' The asymmetry is the same one section~6 analyzed, seen from the side of the good rather than of being: creatures are lovable to God because he made them so, not the reverse.} We love what is already worth loving; that love is a response, and the beloved's goodness precedes it. Divine love cannot be a response, because there was nothing there to respond to. It is the reason anything is there to be responded to at all.

Two texts deserve to be set beside that conclusion. The first is the sentence Aquinas cites as his scriptural warrant, from the Book of Wisdom: ``Thou lovest all things that are, and hatest none of the things which Thou hast made.''\endnote{\emph{Summa theologiae} I, q.~20, a.~2, \latin{sed contra}, quoting Wisdom 11:25 in the Douay numbering, in the Dominican translation's rendering; checked 2026-07-25 as above. Scripture is not independently quoted or exegeted anywhere in this essay; it appears here and in the following note inside the wording of the two witnesses being cited, in their own renderings.} The second is the Catechism's, and its placement is the striking thing: having stated that God upholds and sustains his creatures at every moment, and that recognizing this dependence is ``a source of wisdom and freedom, of joy and confidence,'' the paragraph closes by quoting the same passage of Wisdom---``For you love all things that exist, and detest none of the things that you have made\ldots{} How would anything have endured, if you had not willed it?''\endnote{\emph{Catechism of the Catholic Church} 301, official Vatican English web text, wording checked 2026-07-25; the paragraph quotes Wisdom 11:24--26 in the rendering the Catechism's English edition uses. That the Church's catechetical statement of conservation and Aquinas's philosophical argument for divine love converge on the same verse is a fact about the two texts, noted here; the essay draws no inference from the convergence beyond what each text states on its own authority.} The same verse stands behind the Church's teaching on conservation and behind the philosophical judgment that God loves what he sustains, because in both cases the reasoning runs the same way: what is held in being is held on purpose, and holding on purpose is what willing good to something means.

To love is to will good to something; to exist at all is already to be receiving good---being, the good under every other good. The council's phrase, quoted earlier as dogma, can now be heard as the exact converse of the vertigo: creation was \emph{not for increasing his own beatitude}. What made the arrest dizzying was that my existence makes no difference to God. What consent discovers inside the same fact is that my existence was therefore never extracted, owed, or useful---that it can only have been \emph{given}. The asymmetry that looked like an abyss is the signature of gift. No one holds a debtor in being gratuitously, at every instant, innermostly, forever, for nothing. The falling was never falling. It was being carried.

One correction to the record belongs here, in fairness to the tradition being used. Section~5 marked a seam: the freedom of creation is dogma, and this essay's compressed ascent did not establish it. That remains true of \emph{this essay}. It should not be taken to mean that the tradition regarded divine freedom as inaccessible to argument. Aquinas argues it explicitly: since God's goodness is perfect and complete without other things, ``inasmuch as no perfection can accrue to Him from them, it follows that His willing things apart from Himself is not absolutely necessary.''\endnote{\emph{Summa theologiae} I, q.~19, a.~3, corpus; checked 2026-07-25, English at New Advent. Aquinas distinguishes what God wills necessarily---his own goodness, the proper object of the divine will---from what he wills of things other than himself, which are willed as ordered to that goodness and are not necessary to it. The distinction between absolute necessity and necessity by supposition is drawn in the same corpus. This essay does not execute or assess the argument; it records that the seam it marked is a seam in its own presentation, not a boundary of the tradition's claims.} The argument turns on exactly the fact this essay has been circling: that nothing is added to God by there being a world. This essay does not run that argument, and its conclusion does not rest on it. But a reader should not be left with the impression that a boundary of these pages is a boundary of reason.

\subsection{Two restraints}

Two restraints keep the conclusion from saying more than it knows.

First, ``loves'' is an analogical word. The judgment that God loves all things is a judgment about what God does---wills and causes good---reached by argument; it is not a window into what loving is \emph{for God}, and natural reason cannot make it one. It does not license the inference that God is fond of things, or moved by them, or affected in the way a lover is affected; the whole of section~6 stands against that. What it licenses is the statement that the good of every creature is willed, and that the willing is not a response to anything outside God.

Even the boldest formulation the tradition permits itself here is carefully fenced. Aquinas quotes Dionysius saying that ``He Himself, the cause of all things, by His abounding love and goodness, is placed outside Himself by His providence for all existing things''---and immediately supplies the sense in which alone it is true: a lover is placed outside himself ``inasmuch as he wills good to the beloved; and works for that good by his providence even as he works for his own.''\endnote{\emph{Summa theologiae} I, q.~20, a.~2, ad~1, quoting Pseudo-Dionysius, \emph{Divine Names} IV.1; checked 2026-07-25, English at New Advent. The objection being answered was that love takes the lover out of himself, which cannot be said of God; Aquinas neither denies the Dionysian sentence nor reads it as a claim about divine change, but glosses ``outside himself'' as willing and providing the good of another. The Dionysian text is quoted here through Aquinas, as he quotes it, and is not separately collated.} The ecstatic language is admitted and then given a determinate meaning that leaves the metaphysics of section~6 intact. That is the discipline of analogy in one gesture: the strong word is kept, and what it can mean here is stated.

Second, the inner life of God is closed to the ascent. Aquinas is categorical that ``it is impossible to attain to the knowledge of the Trinity by natural reason,''\endnote{\emph{Summa theologiae} I, q.~32, a.~1, corpus; checked 2026-07-25 as above. The article adds the reason---natural reason reaches God through creatures, and creatures manifest the divine creative power which is common to the whole Trinity, not the distinction of persons---and the warning that attempting such proofs exposes faith to derision. That God is triune, and that the love reached analogically by reason is, in its own life, Father, Son, and Spirit, are truths of revelation; this essay's argument neither reaches nor gestures at them, and the series to which it belongs takes them up only where revelation is explicitly on the table.} and the essay's philosophical thread ends there. What has been reached is that the source wills the good of everything it sustains. What that willing is, in itself, has not been reached and will not be.

\subsection{The failure of the words, which is not a failure of the argument}

There is a last discipline, older than the scholastic one and harder to hold. Even the surviving words must be handled as the tradition's masters handled them---as true, and as breaking in the hand.

The anonymous author who wrote as Dionysius, at the border of Greek philosophy and Christian prayer, counseled the mind ascending past every concept toward ``Him Who is above every essence and knowledge.'' His method is not the accumulation of better descriptions but their removal: we celebrate the source ``through the abstraction of all existing things, just as those who make a lifelike statue, by extracting all the encumbrances which have been placed upon the clear view of the concealed, and by bringing to light, by the mere cutting away, the genuine beauty concealed in it.''\endnote{Pseudo-Dionysius, \emph{Mystical Theology}, ch.~1 (the ascent ``to the union, as attainable, with Him Who is above every essence and knowledge'') and ch.~2 (the statue and the cutting away); quoted from the public-domain Parker translation (1897), wording checked 2026-07-25 at the CCEL transcription. The author is an anonymous Christian writer of the late fifth or early sixth century, conventionally cited under his pseudonym; he is not the Athenian convert of Acts 17.} And so the denials at the end: the source ``is neither soul, nor mind, nor has imagination, or opinion, or reason, or conception; neither is expressed, nor conceived,'' and neither ``lives, nor is life; neither is essence nor eternity''---not because it is less than these, but because it is not one of the things that \emph{have} them.\endnote{\emph{Mystical Theology}, ch.~5, Parker translation, checked 2026-07-25 as above. Cited for apophatic restraint only: the denials mark the failure of creaturely categories at the source of categories, not a retraction of the analogical judgments made above. The two moves are compatible in the tradition that made both, and the appendix records that this essay does not undertake the technical reconciliation of analogy and negation.}

Apophasis of that kind is not agnosticism, and it is certainly not the discovery that the argument failed. It is precision about a source the argument itself showed to be beyond the genus of borrowed things: exactly what one should expect to find at the end of a road whose first step was that existence is not a feature. The creature's best sentences about God are true the way a map is true of a country: reliably, and without resemblance.

\subsection{The fork, and what lies past it}

So Part I closes at the fork it uncovered. A mind that follows the argument to its end stands arrested: wholly real, genuinely acting, and not for one instant self-possessed---held, innermostly, by what it cannot repay and was never asked to. Recoil reads the holding as humiliation and is wrong about the picture, not the scale. Consent reads it as gift, has the better of the argument, and finds that its natural terminus is not a proposition but a posture---the reverence the tradition calls adoration, owed to the excellence that gives and does not receive.\endnote{On adoration as the act of reverence rendered to God on account of the divine excellence, see \emph{Summa theologiae} II-II, q.~84, a.~1, corpus (``Adoration is directed to the reverence of the person adored\ldots{} it is proper to religion to show reverence to God'') and ad~1 (``Reverence is due to God on account of His excellence''); checked 2026-07-25, English at New Advent. The step from consent to adoration is here only named, not developed; Part II of this series takes it up, including the bodily acts treated in the same question at a.~2, which this essay does not enter.}

What that posture looks like when an embodied life---a creature of mornings, obligations, and a body---takes it up is the matter of Part II; what the source has revealed about itself beyond the ascent's reach, and what it makes of creatures who consent, is the matter of Part III. The present essay claims what it argued: that the vertigo is the truth; that the truth was never a threat; and that the ground beneath the falling creature is not ground but hands, and always was.
